\begin{proof}[Proof of \cref{obj:lem:factorial-moment-mean}]
Write
\[
\mu_\infty=P^{\otimes\infty},\qquad
M=M(n),\qquad B=B(n),\qquad m_1=|I_1|,
\]
where \(\mu_\infty\) is the ambient iid product law.  For a cell
\(c=(a,y)\in\{0,1\}^2\), write
\[
q_c=q_{aky}=P(X=k,A=a,Y=y),
\qquad
q_k=(q_c)_{c\in\{0,1\}^2}.
\]
Since \(n\ge N_0\), the calibration in \cref{obj:def:hybrid-estimator-handle} gives
\[
4M^2\le m_1 .
\]
Because \(M\ge2\), this implies \(M\le m_1\).
% lean: CausalSmith.Stat.DiscreteAteMinimaxLoggap.calibrationCutoff

For any multi-index \(r=(r_c)_c\) with \(|r|\le m_1\), the normalized factorial monomial satisfies
\[
\int
\mathfrak m_{k,r}\,d\mu_\infty
=
\prod_{c\in\{0,1\}^2} q_c^{\,r_c}.
\]
Indeed, on the estimation split \(I_1\), define
\[
L_i=
\begin{cases}
(A_i,Y_i),& X_i=k,\\
\star,& X_i\ne k .
\end{cases}
\]
The variables \(L_i\), \(i\in I_1\), are iid, and \(P(L_i=c)=q_c\).  Expanding
\(\prod_c (N^{(1)}_{kc})_{r_c}\) counts ordered injective assignments of \(|r|\) requested labels to distinct indices in \(I_1\).  There are \((m_1)_{|r|}\) such assignments, each has probability \(\prod_c q_c^{r_c}\), and division by \((m_1)_{|r|}\) gives the displayed identity.
% lean: CausalSmith.Stat.DiscreteAteMinimaxLoggap.integral_factorialMonomial_trunc

For the sparse arms in \cref{obj:lem:sparse-arm-expansion}, write
\[
\widehat A_{a,k}
=
\sum_{j=0}^{M-2}\sum_{t=0}^{j}\sum_{c\in\{0,1\}^2}
B^{-1}g_{M,j}B^{-j}\binom jt\,
\mathfrak m_{k,r^{a,c,j,t}} .
\]
That lemma gives the pointwise identity
\[
\widehat P_k=\widehat A_{1,k}-\widehat A_{0,k}.
\]
% lean: CausalSmith.Stat.DiscreteAteMinimaxLoggap.factorialPolynomialContribution_eq_sparseArms

For \(a\in\{0,1\}\), \(j\in\{0,\ldots,M-2\}\), \(t\in\{0,\ldots,j\}\), and \(c\in\{0,1\}^2\), the sparse exponent
\[
r^{a,c,j,t}
=
e_c+e_{a1}+t\,e_{a0}+(j-t)\,e_{a1}
\]
has total degree
\[
|r^{a,c,j,t}|=j+2\le M\le m_1 .
\]
% lean: CausalSmith.Stat.DiscreteAteMinimaxLoggap.multiDegree_factorialExpansionIndex
Therefore the factorial-monomial identity applies in every summand of each sparse arm.  Finite-sum linearity gives, for each \(a\in\{0,1\}\),
\[
\int \widehat A_{a,k}\,d\mu_\infty
=
\sum_{j=0}^{M-2}\sum_{t=0}^{j}\sum_{c\in\{0,1\}^2}
B^{-1}g_{M,j}B^{-j}\binom jt\,
q_k^{\,r^{a,c,j,t}} .
\]
% lean: CausalSmith.Stat.DiscreteAteMinimaxLoggap.integral_factorialPolynomialContribution_trunc

By the definition of \(r^{a,c,j,t}\),
\[
q_k^{\,r^{a,c,j,t}}
=
q_c\,q_{a1}\,q_{a0}^{\,t}\,q_{a1}^{\,j-t}.
\]
Thus, for each fixed \(a\) and \(j\),
\[
\sum_{t=0}^{j}\sum_{c\in\{0,1\}^2}
\binom jt q_k^{\,r^{a,c,j,t}}
=
\left(\sum_{c\in\{0,1\}^2}q_c\right)
q_{a1}
(q_{a0}+q_{a1})^j,
\]
by the binomial theorem.
% lean: CausalSmith.Stat.DiscreteAteMinimaxLoggap.factorialExpansionIndex_binomial_sum
Consequently, for each \(a\in\{0,1\}\),
\[
\int \widehat A_{a,k}\,d\mu_\infty
=
\frac{s(q_k)t_a(q_k)}{B}
\sum_{j=0}^{M-2}
g_{M,j}
\left(\frac{s_a(q_k)}{B}\right)^j
=
\frac{s(q_k)t_a(q_k)}{B}
G_M\!\left(\frac{s_a(q_k)}{B}\right),
\]
where \(t_1(q_k)=q_{11}\) and \(t_0(q_k)=q_{01}\).
% lean: CausalSmith.Stat.DiscreteAteMinimaxLoggap.integral_factorialPolynomialContribution_trunc

Subtracting the two arm identities and using \cref{obj:lem:sparse-arm-expansion},
\[
\int \widehat P_k\,d\mu_\infty
=
\frac{s(q_k)t_1(q_k)}{B}
G_M\!\left(\frac{s_1(q_k)}{B}\right)
-
\frac{s(q_k)t_0(q_k)}{B}
G_M\!\left(\frac{s_0(q_k)}{B}\right)
=
P_{M,B}(q_k),
\]
where the final equality is the definition of the light-cell polynomial in \cref{obj:def:hybrid-estimator-handle}.  Substituting \(M=M(n)\) and \(B=B(n)\) gives the claimed identity.
% lean: CausalSmith.Stat.DiscreteAteMinimaxLoggap.integral_factorialPolynomialContribution_trunc
\end{proof}
